\subsection{Proof of Correctness}

Throughout, for $lo \le hi$ let
$\mathcal{I}(lo, hi) = \{[l, r] : lo \le l \le r \le hi\}$
denote the set of periods inside $[lo, hi]$.

\medskip
\noindent\textbf{Lemma 1 (Crossing Periods).}
Let $lo \le mid < hi$. Then
\textsc{Max-Crossing}$(X, lo, mid, hi)$ returns
$(l', r', s')$ with $lo \le l' \le mid < r' \le hi$,
$s' = S(l', r')$, and

\[
s' = \max \{ S(l, r) : lo \le l \le mid < r \le hi \}.
\]

\noindent\textbf{Proof.}
Any crossing period $[l, r]$ splits into two
non-empty parts, $[l, mid]$ and $[mid+1, r]$, so

\[
S(l, r) = S(l, mid) + S(mid+1, r).
\]

The first term depends only on $l$ and the second
only on $r$, and every pair $(l, r)$ with
$lo \le l \le mid$ and $mid+1 \le r \le hi$ is a
valid crossing period. Hence the two terms can be
maximized independently:

\[
\max_{l, r} S(l, r)
= \max_{lo \le l \le mid} S(l, mid)
+ \max_{mid < r \le hi} S(mid+1, r).
\]

In the first loop of \textsc{Max-Crossing}, the
following invariant holds after processing index
$i$: $sum = S(i, mid)$, and $leftSum = S(bestL, mid)$
is the maximum of $S(l, mid)$ over $i \le l \le mid$.
It holds after the first iteration ($i = mid$)
because $leftSum$ starts at $-\infty$, and each later
iteration adds $X[i]$ to $sum$ and updates $leftSum$
only if the new suffix is larger. When the loop
ends at $i = lo$, $leftSum$ is the maximum over all
$lo \le l \le mid$. A symmetric invariant shows that
$rightSum$ is the maximum of $S(mid+1, r)$ over
$mid < r \le hi$. Therefore $s' = leftSum + rightSum$
is the maximum crossing sum, attained by
$[bestL, bestR]$. \hfill $\square$

\medskip
\noindent\textbf{Theorem (Correctness of Max-Subarray).}
For every $lo \le hi$,
\textsc{Max-Subarray}$(X, lo, hi)$ returns
$(l, r, s)$ with $[l, r] \in \mathcal{I}(lo, hi)$,
$s = S(l, r)$, and

\[
s = \max_{[l', r'] \in \mathcal{I}(lo, hi)} S(l', r').
\]

\noindent\textbf{Proof.}
By strong induction on the length
$m = hi - lo + 1$.

\emph{Base case} ($m = 1$). Then
$\mathcal{I}(lo, lo) = \{[lo, lo]\}$, and the
algorithm returns $(lo, lo, X[lo])$, which is the
only and therefore the best period.

\emph{Inductive step} ($m \ge 2$). Assume the
claim holds for all lengths less than $m$.
Since $lo < hi$, $mid = \lfloor (lo+hi)/2 \rfloor$
satisfies $lo \le mid < hi$, so both halves
$[lo, mid]$ and $[mid+1, hi]$ are non-empty and
have length strictly less than $m$.

The set of periods partitions as

\[
\mathcal{I}(lo, hi) = \mathcal{I}(lo, mid)
\,\cup\, \mathcal{I}(mid+1, hi)
\,\cup\, \mathcal{C},
\]

where $\mathcal{C} = \{[l, r] : lo \le l \le mid < r \le hi\}$
is the set of crossing periods. This holds because,
for any $[l, r] \in \mathcal{I}(lo, hi)$, either
$r \le mid$ (left), or $l > mid$ (right), or
$l \le mid < r$ (crossing).

By the inductive hypothesis, $L$ and $R$ are optimal
over $\mathcal{I}(lo, mid)$ and $\mathcal{I}(mid+1, hi)$,
respectively. By Lemma 1, $C$ is optimal over
$\mathcal{C}$. The maximum over a union of sets equals
the maximum of the maxima over each set, so the
candidate among $L, R, C$ with the largest sum is
optimal over $\mathcal{I}(lo, hi)$. Each of $L, R, C$
is a valid period in $\mathcal{I}(lo, hi)$ whose
reported sum equals its true sum, so the returned
triple is also valid. \hfill $\square$

\medskip
Applying the theorem with $lo = 1$ and $hi = n$
shows that the algorithm returns a period of maximum
total queue build-up over the whole observation
horizon. The algorithm terminates because each
recursive call is on a strictly shorter range.
